\documentclass{article}

\usepackage{amsmath}
\usepackage{amssymb}

\title{CSE400 -- Fundamentals of Probability in Computing}
\author{Scribe: Joint Random Variables and Joint Distributions}
\date{Instructor: Dhaval Patel \\ SEAS -- Ahmedabad University \\ March 10, 2026}

\begin{document}

\maketitle

\section*{1. Joint Random Variables}

In many real-world applications, we are interested in \textbf{more than one random variable simultaneously}.

Examples include:

\begin{itemize}
\item Smart transportation systems
\item Wireless network performance
\item Weather prediction
\item Drone delivery systems
\item Rolling two dice
\end{itemize}

These situations require studying \textbf{joint random variables}.

\section*{2. Motivation for Studying Joint Random Variables}

\subsection*{2.1 Smart Transportation}

Two random variables are considered:

\begin{itemize}
\item $X$ = Vehicle Speed
\item $Y$ = Traffic Density
\end{itemize}

Question:

Can speed be high when traffic density is also high?

This motivates the study of \textbf{correlated random variables}.

Applications include:

\begin{itemize}
\item Adaptive traffic signals
\item Congestion prediction
\item Autonomous vehicle routing
\end{itemize}

\subsection*{2.2 Wireless Network Performance}

Two random variables are considered:

\begin{itemize}
\item $X$ = Signal-to-Noise Ratio (SNR)
\item $Y$ = Data Throughput
\end{itemize}

Question:

If SNR is high, will throughput always be high?

Reason: Throughput depends on SNR and also varies due to:

\begin{itemize}
\item Fading
\item Interference
\item Network congestion
\end{itemize}

Applications include:

\begin{itemize}
\item 5G scheduling
\item Adaptive modulation
\item Network planning
\end{itemize}

\subsection*{2.3 Other Examples}

Examples of joint random variables include:

\begin{itemize}
\item Weather prediction
\begin{itemize}
\item $X$ = Rainfall
\item $Y$ = Temperature
\end{itemize}

\item Drone delivery
\begin{itemize}
\item $X$ = Wind Speed
\item $Y$ = Delivery Time
\end{itemize}

\item Rolling two dice
\begin{itemize}
\item $X$ = Outcome of first die
\item $Y$ = Outcome of second die
\end{itemize}
\end{itemize}

Events such as:

\[
X + Y = 7
\]

or

\[
X > Y
\]

require the use of \textbf{joint probability mass functions}.

\section*{3. Discrete Case -- Joint PMF}

\subsection*{Definition}

For discrete random variables $X$ and $Y$, the \textbf{joint probability mass function (PMF)} is defined as

\[
p_{X,Y}(x,y) = P(X = x, Y = y)
\]

This gives the probability that:

\begin{itemize}
\item $X = x$
\item $Y = y$
\end{itemize}

simultaneously.

\subsection*{Joint PMF Table Representation}

For discrete random variables, the joint PMF can be represented using a \textbf{table}, where:

\begin{itemize}
\item Rows correspond to values of $X$
\item Columns correspond to values of $Y$
\item Each cell contains $P(X=x, Y=y)$
\end{itemize}

\section*{4. Example: Rolling Two Dice}

Let:

\begin{itemize}
\item $X$ = outcome of the first die
\item $Y$ = outcome of the second die
\end{itemize}

Possible values:

\[
X \in \{1,2,3,4,5,6\}
\]

\[
Y \in \{1,2,3,4,5,6\}
\]

Each pair $(X,Y)$ corresponds to one outcome of the experiment.

The joint PMF assigns probability:

\[
P(X=x, Y=y) = \frac{1}{36}
\]

for each possible pair $(x,y)$.

The joint PMF can therefore be represented using a \textbf{6×6 probability table}.

\section*{5. Continuous Case -- Joint PDF}

For continuous random variables, we use the \textbf{joint probability density function (PDF)}.

\subsection*{Definition}

The \textbf{joint PDF} of continuous random variables $X$ and $Y$ is denoted as:

\[
f_{X,Y}(x,y)
\]

and satisfies:

\[
P((X,Y) \in A) = \int\int_A f_{X,Y}(x,y)\,dx\,dy
\]

for a region $A$.

\subsection*{Geometric Interpretation}

For continuous joint distributions:

\begin{itemize}
\item The joint PDF defines a \textbf{surface over the $(x,y)$ plane}.
\item Probabilities correspond to \textbf{volumes under this surface}.
\end{itemize}

\section*{6. Worked Example}

A joint PDF $f_{X,Y}(x,y)$ is given.

The example requires computing probabilities using integration over a region.

The solution proceeds through the following steps:

\begin{enumerate}
\item Identify the \textbf{valid region of support} for $X$ and $Y$
\item Set up the \textbf{double integral} over the region
\item Evaluate the integral step by step
\item Verify that the PDF integrates to \textbf{1 over the support}
\end{enumerate}

The detailed calculations follow the steps shown in the material.

\section*{7. Joint Cumulative Distribution Function (Joint CDF)}

\subsection*{Definition}

The \textbf{joint cumulative distribution function (CDF)} is defined as

\[
F_{X,Y}(x,y) = P(X \le x, Y \le y)
\]

This represents the probability that:

\begin{itemize}
\item $X$ is less than or equal to $x$, and
\item $Y$ is less than or equal to $y$.
\end{itemize}

\section*{8. Joint CDF -- Continuous Case}

For continuous random variables:

\[
F_{X,Y}(x,y)
=
\int_{-\infty}^{x}
\int_{-\infty}^{y}
f_{X,Y}(u,v)\, dv\, du
\]

This integrates the joint PDF over the region:

\[
(-\infty, x] \times (-\infty, y]
\]

\section*{9. Joint CDF -- Discrete Case}

For discrete random variables:

\[
F_{X,Y}(x,y)
=
\sum_{u \le x}
\sum_{v \le y}
P(X=u, Y=v)
\]

The joint CDF is obtained by summing the probabilities of all pairs $(u,v)$ such that:

\[
u \le x, \quad v \le y
\]

\section*{10. Properties of the Joint CDF}

The joint CDF satisfies the following properties.

\subsection*{Property 1}

\[
0 \le F_{X,Y}(x,y) \le 1
\]

\subsection*{Property 2}

The function is \textbf{non-decreasing} in both variables.

\subsection*{Property 3}

Limits:

\[
F_{X,Y}(-\infty,-\infty) = 0
\]

\[
F_{X,Y}(\infty,\infty) = 1
\]

\subsection*{Property 4}

Probabilities over rectangular regions can be computed using the joint CDF.

\section*{11. Summary}

The lecture introduces the concept of \textbf{joint random variables} and their distributions.

Topics covered include:

\begin{itemize}
\item Motivation for studying multiple random variables
\item Joint PMF for discrete variables
\item Joint PDF for continuous variables
\item Joint CDF definitions
\item Properties of joint CDF
\end{itemize}

These concepts allow the analysis of systems where \textbf{multiple random variables interact simultaneously}.

\end{document}